%! Radicals and Rational Exponents — Unit 1, Lesson 1.2 — Lesson Plan
\documentclass[10pt]{article}
\usepackage{atda-article}
\usepackage{atda-boxes}
\usepackage{graphicx}   % -article does NOT load graphicx; needed for thumbnails/screenshots
\graphicspath{{images/}}
\newcommand{\TallMath}[1]{$\displaystyle #1\rule[-1.4em]{0pt}{3.2em}$}
\newcommand{\UnitNumberName}{Unit 1: Algebraic Expressions and Factoring \quad}
\newcommand{\LessonNumberName}{Lesson 1.2: Radicals and Rational Exponents}

\begin{document}
\begin{center}
    {\Large\bfseries \CourseName \\
    {\small\bfseries \UnitNumberName \LessonNumberName}}
\end{center}

\begin{tcolorbox}[colback=mist, colframe=royal, boxrule=0.9pt, arc=2mm,
    left=2mm, right=2mm, top=1.5mm, bottom=1.5mm]
    \textbf{Primary Objective:} I can simplify a radical with a numeric or algebraic radicand by
    pulling out every perfect $n$th-power factor, and I can add, subtract, multiply, and divide
    radicals --- including clearing a radical out of a denominator. I can move between a radical
    and a rational exponent in both directions and evaluate expressions like $27^{2/3}$ without a
    calculator. I can use a square-root model from a real situation --- speed from a skid mark,
    fall time from a tower height --- to answer a question and say what the answer means.
\end{tcolorbox}

\renewcommand{\arraystretch}{1.6}
\begin{skillbox}[Priority Ideas \& Skills]{goldbox}
\begin{tabularx}{\linewidth}{@{}X|X@{}}
\begin{minipage}[t]{\linewidth}\vspace{0pt}
\textbf{Knowledge \& Skills covered --- \texttt{A2.EO.2a--c}:}
\begin{itemize}[leftmargin=1.1em, itemsep=2pt, topsep=2pt]
  \item \texttt{2a} --- Simplify and determine equivalent radical expressions with
        \textbf{numeric and algebraic radicands}.
  \item \texttt{2b} --- \textbf{Add, subtract, multiply, and divide} radical expressions,
        simplifying the result; includes \textbf{rationalizing the denominator}.
  \item \texttt{2c} --- \textbf{Convert} between radical expressions and expressions containing
        \textbf{rational exponents}, in both directions.
\end{itemize}
\textbf{Supporting fluency (Lesson 1.1, not a new code):}
\begin{itemize}[leftmargin=1.1em, itemsep=2pt, topsep=2pt]
  \item The five exponent rules and the zero/negative exponents --- all of them survive
        unchanged when the exponent becomes a fraction.
  \item Like terms and ``every term times every term'' --- both reappear verbatim, with
        $\sqrt3$ in the role of $x$.
\end{itemize}
\textbf{Where these codes go next:} \texttt{A2.EO.3b} in Lessons 1.3--1.5 (pulling out a common
factor is the same habit as pulling out a perfect square); \texttt{A2.EO.1} in 1.7, where
rationalizing becomes a rational-expression move; \texttt{A2.F.1} in Unit 3, where $a^{m/n}$ is
what makes an exponential function defined between the integers.
\end{minipage}
&
\begin{minipage}[t]{\linewidth}\vspace{0pt}
\textbf{Key Understandings (the \emph{why}):}
\begin{itemize}[leftmargin=1.1em, itemsep=2pt, topsep=2pt]
  \item \textbf{A radical undoes a power, so it scales the opposite way.} Lesson 1.1's hook said
        doubling a radius quadruples the area. Today four times the skid length gives only twice
        the speed. Same machinery, run backwards.
  \item \textbf{``Simplified'' means the \emph{largest} perfect power came out.}
        $\sqrt{72}=\sqrt4\cdot\sqrt{18}$ is a true statement and a wrong answer --- the usual
        cause is grabbing the first square factor instead of the biggest.
  \item \textbf{Rational exponents are forced, not defined.} $9^{1/2}\cdot 9^{1/2}=9$ by the
        product rule, so $9^{1/2}$ \emph{has} to be $\sqrt9$. Presenting $a^{m/n}$ as a
        definition to memorize is what produces $16^{1/2}=8$.
  \item \textbf{Like radicals are like terms.} $\sqrt3$ behaves exactly like $x$ --- which is why
        the instruction is always \emph{simplify first, then decide what combines}.
  \item \textbf{A root splits over a product, never over a sum.} $\sqrt{9+16}=5$, not $7$. This
        is the single most durable error in the lesson and it returns in Unit 7 the moment
        students meet the Pythagorean Theorem.
  \item \textbf{Exact beats rounded.} $30\sqrt2$ carries more information than $42.4$; the
        decimal is for the sentence at the end, not for the work in the middle.
\end{itemize}
\end{minipage}
\end{tabularx}
\end{skillbox}

\begin{skillbox}[{Vocabulary, Concepts, \& Theorems}]{greenbox}
\renewcommand{\arraystretch}{1.35}
\begin{tabularx}{\linewidth}{@{}p{3.6cm}X@{}}
\textbf{Term} & \textbf{Meaning students record} \\
\hline
Radical, radicand, index & In $\sqrt[n]{a}$ the whole symbol is the radical, $a$ is the radicand,
    and $n$ is the index. A missing index means $2$. \\
Principal square root & The non-negative root: $\sqrt{49}=7$. The other one is written
    $-\sqrt{49}$. \\
Product \& quotient properties & $\sqrt[n]{ab}=\sqrt[n]{a}\sqrt[n]{b}$ and
    $\sqrt[n]{a/b}=\sqrt[n]{a}/\sqrt[n]{b}$ for $b \ne 0$. There is \emph{no} such property for
    sums. \\
Simplified radical form & No perfect $n$th-power factor remains inside, and no radical remains in
    a denominator. \\
Like radicals & Same index \emph{and} same radicand --- the only radicals that may be added or
    subtracted. \\
Rational exponent & $a^{m/n}=\sqrt[n]{a^m}=\left(\sqrt[n]{a}\right)^m$. Denominator $=$ index,
    numerator $=$ power. \\
Rationalizing the denominator & Multiplying by a form of $1$ to clear a radical downstairs:
    $\dfrac{6}{\sqrt3}\cdot\dfrac{\sqrt3}{\sqrt3}=2\sqrt3$. \\
Conjugate & $\sqrt a + b$ and $\sqrt a - b$. Their product is $a-b^2$ --- no radical --- which is
    why it clears a two-term denominator (Tier E). \\
\end{tabularx}
\end{skillbox}

\begin{fixedskillbox}[Activate Prior Knowledge \& Spiral Review]{mist}
\begin{tabularx}{\linewidth}{@{}X c@{}}
\begin{minipage}[t]{0.55\linewidth}\vspace{0pt}
\textbf{Four items, about 6 minutes:}
\begin{enumerate}[leftmargin=1.4em, itemsep=1pt, topsep=2pt]
  \item Spiral from 1.1: simplify $\left(5y^3\right)^2$
  \item Evaluate $\sqrt{144}$, $\sqrt[3]{64}$, $\sqrt{81}$ (perfect-power recall)
  \item Rewrite $72$, $48$, $200$ as (largest perfect square) $\times$ (the rest)
  \item \emph{Discovery}: apply the product rule to $9^{1/2}\cdot 9^{1/2}$
\end{enumerate}
\end{minipage}
&
\begin{minipage}[t]{0.38\linewidth}\vspace{0pt}
\textbf{How to use the results:} item 3 is the whole simplification algorithm in disguise --- if a
student writes $72=4\cdot 18$, that is exactly the error that will cost them box 1, so catch it
here. Item 4 is not review at all; it is the lesson's first two minutes. Take it out loud
\emph{before} the hook and let a student say ``so $9^{1/2}$ is $3$.''
\end{minipage}
\end{tabularx}
\end{fixedskillbox}

\begin{skillbox}[Hook]{mist}
\textbf{Reading a skid mark.} Crash investigators estimate a car's speed from the length of its
skid marks: on a road with drag factor $f$, a skid of $d$ feet means the car was going about
$s=\sqrt{30fd}$ mph. Two cars skid to a stop on the same dry asphalt ($f=0.8$): Car A leaves a
$60$-foot skid, Car B a $240$-foot skid --- \textbf{four times as long}. \textbf{Pose it this
way:} ``Four times the skid. Was Car B going four times as fast? Vote before you compute.'' Most
rooms vote yes. Then compute:
\[ \sqrt{30(0.8)(60)} = \sqrt{1440} = 12\sqrt{10} \approx 37.9 \qquad
   \sqrt{30(0.8)(240)} = \sqrt{5760} = 24\sqrt{10} \approx 75.9. \]
Four times the skid, exactly \emph{twice} the speed, because $\sqrt4=2$. \textbf{Name the
connection to Lesson 1.1 out loud:} yesterday doubling a radius quadrupled an area; today
quadrupling a distance only doubles a speed. A radical is that machinery running backwards --- and
$\sqrt{1440}=12\sqrt{10}$ is the first reason simplifying a radical is worth doing.
\end{skillbox}

\boxguard
\begin{skillbox}[Lesson]{mist}
\begin{multicols}{2}
\textbf{1. Simplifying (12 min).} Start from warm-up item 3 --- the largest perfect square, not
the first one. State both properties, then fill the table together. \textbf{Ask:} ``Why is
$\sqrt4\cdot\sqrt{18}$ a wrong answer even though it is a true statement?'' Do the algebraic
radicands slowly: $\sqrt{72x^5}$ works because $x^5=x^4\cdot x$, and students who see the even
exponent as the ``perfect square part'' generalize correctly to $\sqrt[3]{40y^7}$. Close the box
by simplifying the hook's $\sqrt{1440}$.

\textbf{2. Rational exponents (12 min).} Run warm-up item 4 back on the board. The point is that
nobody chose $a^{1/n}=\sqrt[n]{a}$ --- the product rule forces it, exactly as the quotient rule
forced $x^0=1$ in Lesson 1.1. \textbf{Ask:} ``Which part of $\tfrac23$ is the root and which is
the power?'' Then $27^{2/3}$: take the root \emph{first} and the arithmetic stays at $3^2$
instead of $\sqrt[3]{729}$. Finish with $\left(32x^{10}\right)^{3/5}$ to show every 1.1 rule
survived intact.

\textbf{3. Operations (12 min).} Frame adding as the like-terms rule with $\sqrt3$ in place of
$x$; that one sentence does most of the work. \textbf{Ask:} ``Are $\sqrt{12}$ and $\sqrt{75}$
like radicals?'' --- they look unlike and are not, which is why \emph{simplify first} is the
standing instruction. Multiplication is 1.1's ``every term times every term'' unchanged. For
rationalizing, insist on the phrase ``multiplying by $1$'' so nobody believes the value changed.

\textbf{4. Guided practice (8 min).} The accident report. Item 1 is the computation, item 2 is
the interpretation (a number \emph{and} a sentence), item 3 is the justification --- and item 3
is the one to protect if time runs short, because it is the hook's idea stated in the students'
own words.
\end{multicols}
\end{skillbox}

\boxguard
\begin{skillbox}[Explicit Instruction: Simplifying Any Radical]{mist}
\begin{tabularx}{\linewidth}{@{}X|X@{}}
\begin{minipage}[t]{\linewidth}\vspace{0pt}
\textbf{The four steps students record:}
\begin{enumerate}[leftmargin=1.4em, itemsep=2pt, topsep=2pt]
  \item Read the index. It tells you which perfect powers to hunt: squares for $n=2$, cubes for
        $n=3$.
  \item Find the \textbf{largest} perfect $n$th-power factor of the radicand --- numbers and
        variables both. For a variable, that is the largest multiple of $n$ in its exponent.
  \item Split the radical over that product and take the root of the perfect part.
  \item Check: is anything left inside that is still a perfect $n$th power? If yes, you did not
        take the largest one.
\end{enumerate}
\textbf{Say this out loud:} ``$\sqrt{a+b}$ is not $\sqrt a + \sqrt b$. The property is about
products.''
\end{minipage}
&
\begin{minipage}[t]{\linewidth}\vspace{0pt}
\textbf{Worked example (algebraic radicand):}
\[
\begin{aligned}
\sqrt[3]{40y^7} &= \sqrt[3]{8y^6 \cdot 5y} \\
                &= \sqrt[3]{8y^6}\cdot\sqrt[3]{5y} \\
                &= 2y^2\sqrt[3]{5y}
\end{aligned}
\]
Index $3$, so hunt cubes: $8$ is the largest cube factor of $40$, and $y^6$ is the largest
multiple-of-$3$ power inside $y^7$. \textbf{The error to name before it happens:} pulling out
$y^7$ whole, or writing $2y^2\sqrt[3]{5}\,y$. The leftover $y$ stays \emph{inside} the radical.
\end{minipage}
\end{tabularx}
\end{skillbox}

\begin{skillbox}[Active Monitoring]{redbox}
\begin{itemize}[leftmargin=1.2em, itemsep=2pt, topsep=2pt]
  \item \textbf{Notes, box 1:} watch for $\sqrt{50}=25\sqrt2$ --- the root was never taken --- and
        for $\sqrt{72x^5}$ answered $6x^2\sqrt2\,x$ with the leftover $x$ escaping the radical.
  \item \textbf{Notes, box 2:} watch for $16^{-1/2}=-4$. It is Lesson 1.1's negative-exponent
        misconception wearing a fraction; ask ``did the exponent ever change a sign yesterday?''
  \item \textbf{Notes, box 3:} the giveaway error is $\sqrt{12}+\sqrt{75}$ declared ``cannot be
        combined.'' Point at the radicands and say nothing else.
  \item \textbf{Activity, Tier A item 3:} students will answer $256$ feet. The fall time depends
        on $\sqrt h$, so doubling it takes \emph{four} times the height. Expect this to be the
        sticking point for a third of the room --- it is the hook again, inverted.
  \item \textbf{Cold-call prompts:} ``Why is $\sqrt4\cdot\sqrt{18}$ not an answer?'' ``Is
        $16^{1/2}$ eight?'' ``Which part of $\tfrac32$ is the root?''
\end{itemize}
\end{skillbox}

\boxguard
\begin{skillbox}[Group Work \& Differentiation]{redbox}
\begin{multicols}{3}
\textbf{Tier R --- Remediate}
\begin{itemize}[leftmargin=1em, itemsep=1pt, topsep=2pt]
  \item Five one-line fluency items: a numeric square root, a cube root, an algebraic radicand, a
        sum of like radicals, and one radical-to-exponent conversion.
  \item Support: have the group list the perfect squares to $144$ on scratch paper first.
\end{itemize}

\columnbreak
\textbf{Tier A --- Approaching Proficiency}
\begin{itemize}[leftmargin=1em, itemsep=1pt, topsep=2pt]
  \item Derive $t=\sqrt h/4$ from $h=16t^2$, then evaluate it at $h=128$.
  \item Double the fall time --- the $512$-foot item, and a sentence on what it costs.
  \item A rationalizing item ($12/\sqrt6$) and a binomial product with radicals.
\end{itemize}

\columnbreak
\textbf{Tier E --- Extension}
\begin{itemize}[leftmargin=1em, itemsep=1pt, topsep=2pt]
  \item Critique two wrong answers, $\sqrt{9+16}=7$ and $16^{1/2}=8$, and name each mistake.
  \item Rationalize $6/\left(\sqrt5-2\right)$ with a conjugate, and say why it works.
  \item Justify: eight times the volume of a cube is only twice the edge.
\end{itemize}
\end{multicols}
\end{skillbox}

\begin{skillbox}[Individual Work \& Assessment]{redbox}
\textbf{Exit ticket (3 items, no notes):} (1) simplify $\sqrt{98x^7}$ --- answer $7x^3\sqrt{2x}$;
(2) combine $\sqrt{27}+\sqrt{12}$ --- answer $5\sqrt3$; (3) \textbf{the conceptual check} --- write
$8^{2/3}$ in radical form, evaluate it ($4$), and say in one sentence what the $3$ does and what
the $2$ does.

\textbf{Using the results:} the three items sample the three lettered skills in order --- $2a$,
$2b$, $2c$ --- so a wrong item names the code to reteach, not just a score out of three. Item 1
failing usually means the largest perfect power was not taken; item 2 failing means the student
never simplified before comparing; item 3 failing means the fraction is still being read as a
fraction rather than as root-and-power. Sort the tickets by \emph{which} item is wrong; that sort
is your Tier R grouping for Lesson 1.3.
\end{skillbox}

\boxguard[18]
\begin{skillbox}[Reinforcement \& Extension]{goldbox}
\textbf{Homework} (13 items): nine fluency items covering all three lettered skills --- simplify
(numeric, cube-root, two-variable algebraic), the four operations including a binomial product and
a rationalized quotient, and conversion in both directions plus two evaluations without a
calculator --- then a three-part Kepler's-third-law model ($T=a^{3/2}$), and one look-ahead item.

\textbf{Extension (optional):} a square gate of side $s$ needs a diagonal brace of $s\sqrt2$;
compute it for $6$ and $12$ feet and explain why doubling the side \emph{doubles} the brace here,
even though quadrupling the drop tower's height only doubled the fall time. The distinction is
that $s\sqrt2$ is $s$ times a constant while $\sqrt h$ is not.

\textbf{Preview of Lesson 1.3 (\texttt{A2.EO.3b}):} factoring by GCF and grouping. Homework item
13 makes the link explicit --- pulling the largest perfect square out of a radicand and pulling
the largest common factor out of a polynomial are the same habit applied to different objects.

\textbf{Connections.} Covers \texttt{A2.EO.2a} (simplify numeric and algebraic radicands),
\texttt{A2.EO.2b} (add, subtract, multiply, divide, rationalize), and \texttt{A2.EO.2c} (convert
radicals $\leftrightarrow$ rational exponents). Builds directly on \texttt{A2.EO.3a} (Lesson 1.1);
feeds \texttt{A2.EO.3b} (1.3--1.5), \texttt{A2.EO.1a--b} (1.7, where rationalizing returns as a
rational-expression move), and \texttt{A2.F.1} in Unit 3, where $a^{m/n}$ is what lets an
exponential function have values between the integers.
\end{skillbox}

% ── Teacher notes — one per component, in packet order ────────────────────────
% Teacher-only prose lives HERE, never in a _key: a note is the one block in a
% key with no counterpart in the blank, so it makes the key longer than the
% blank. See references/conventions.md ("teachernote").
\begin{teachernote}[Warm-Up]
    \textbf{6 minutes, and take item 4 out loud before anything else.} Answers: (1) $25y^6$;
    (2) $12$, $4$, $9$; (3) $36\cdot2$, $16\cdot3$, $100\cdot2$; (4) $9^{1/2}\cdot9^{1/2}=9^1=9$,
    so $9^{1/2}=3$. Item 4 is the pivot of the whole lesson --- if you collect it silently you
    have spent the discovery and bought nothing. Put $9^{\frac12+\frac12}$ on the board and ask
    who can finish the sentence. Item 3 is the diagnostic that matters most: a student who writes
    $72=4\cdot18$ will produce $\sqrt{72}=2\sqrt{18}$ within ten minutes, so correct the habit
    now, at the level of whole numbers, where it is cheap.
\end{teachernote}

\begin{teachernote}[Guided Notes]
    \textbf{The tables in boxes 1 and 2 are fill-ins, not handouts.} Answers, box 1:
    $5\sqrt2$, $3\sqrt[3]2$, $6x^2\sqrt{2x}$, $2y^2\sqrt[3]{5y}$, $\dfrac74$; the in-text blank is
    \emph{simplified}. Box 2: $x^{5/3}$, $y^{3/4}$, $t^{5/6}$ down the left, and $9$, $\dfrac14$,
    $27$ down the right; the blank is \emph{index}. Box 3's blank is \emph{radicand}. Box 2 is the
    box that pays off all year --- it is why Unit 3's exponential functions have values between
    the integers --- so do not cut it for time. If the period is running behind, cut the binomial
    product in box 3 (keep the rationalizing example) rather than shortening guided practice. On
    the practice item 3, expect ``$240$ feet''; let the group compute $\sqrt{30(0.5)(240)}$ and
    watch the mismatch send them back.
\end{teachernote}

\begin{teachernote}[Group Activity]
    \textbf{Groups of three, 15 minutes, all groups start at Tier R.} Tier R is five one-liners; a
    group still working on it at six minutes needs you, not more time. Tier A item 1 is worth
    doing on the board first if more than two groups stall --- solving $h=16t^2$ for $t$ is the
    algebra, and the rest of the tier is arithmetic once it is done. Item 3 is the lesson's real
    assessment: doubling the fall time takes $512$ feet, not $256$, and roughly a third of the
    room will write $256$. Do not correct it --- ask them to compute $\sqrt{256}/4$ and compare it
    to $2\sqrt2$. Tier E item 1 is the one to share out; both errors ($\sqrt{9+16}=7$ and
    $16^{1/2}=8$) will be in the room whether or not anyone admits to them.
\end{teachernote}

\begin{teachernote}[Exit Ticket]
    \textbf{5 minutes, notes closed, hand it to me at the door.} The three items are
    \texttt{2a}, \texttt{2b}, \texttt{2c} in order, so read the tickets by \emph{which} item is
    wrong rather than by a score. A student who misses only item 1 is grabbing the first perfect
    square instead of the largest; only item 2 means they compared $\sqrt{27}$ and $\sqrt{12}$
    before simplifying; only item 3 means the fraction is still being read as a fraction. Item 3
    asks for a sentence: accept anything that identifies the $3$ as the root and the $2$ as the
    power, and do not require the words ``index'' or ``radicand.''
\end{teachernote}

\begin{teachernote}[Homework]
    \textbf{Expect 20--25 minutes.} Items 1--9 are deliberately unmixed by type so a student can
    tell you which \emph{kind} of problem broke. Item 3 is the only one with two variables and an
    odd exponent under a square root, so check it first when you spot-grade --- $y^3$ splitting
    into $y^2\cdot y$ is the step most likely to be skipped. Items 10--12 are the first place in
    this course where a rational exponent does real scientific work; item 12 needs a calculator
    and item 11 deliberately does not, so students see that $4^{3/2}$ is exact arithmetic. Item 13
    is the bridge into Lesson 1.3 --- students are asked to pull out a shared factor, framed as
    the same move they spent all period making with perfect squares, so nobody needs vocabulary
    they have not met. The extension is genuinely optional and is the best follow-up for whoever
    finished Tier E early.
\end{teachernote}
\end{document}
